\section{Water Usage}

Water management represents a critical environmental consideration for Indonesian companies, particularly given the country's archipelagic geography and vulnerability to climate change impacts. This section examines water consumption patterns, conservation efforts, and regulatory compliance across major industrial sectors in Indonesia.

\subsection{Industrial Water Consumption Patterns}

Indonesia's industrial water consumption has grown significantly over the past decade, with manufacturing, mining, and energy sectors accounting for approximately 15\% of total freshwater withdrawals \cite{ADB2020}. The pulp and paper industry remains one of the largest industrial water consumers, with major players like Asia Pulp \& Paper (APP) and APRIL reporting annual water usage exceeding 100 million cubic meters per facility \cite{APP2022}.

Textile manufacturing, concentrated in West Java and Central Java, presents unique water management challenges due to dyeing and finishing processes. The Bandung Industrial Area reports average water consumption of 25-30 cubic meters per ton of fabric produced, though water-efficient technologies can reduce this to 15-20 cubic meters \cite{BPS2021}.

Mining operations, particularly coal mining in East Kalimantan and gold mining in Papua, demonstrate varied water usage patterns. Open-pit coal mining requires substantial water for dust suppression and coal washing, with water intensity ranging from 2-8 cubic meters per ton of coal produced depending on geological conditions \cite{MiningIndonesia2023}.

\subsection{Water Stress and Regional Considerations}

Indonesia's water stress varies significantly by region. Java, home to 60\% of the population but containing only 10\% of the nation's water resources, faces severe water stress \cite{WorldResourcesInstitute2021}. Industrial facilities in Greater Jakarta, Surabaya, and Bandung operate in areas classified as high to extremely high water stress.

The government's National Water Resources Policy (Peraturan Pemerintah No. 121/2015) establishes water allocation priorities, with domestic and environmental needs taking precedence over industrial use. This regulatory framework has prompted companies to implement water recycling and conservation measures, particularly in water-stressed regions.

\subsection{Conservation Technologies and Practices}

Leading Indonesian companies have adopted various water conservation technologies. PT Unilever Indonesia's Cikarang factory achieved a 45\% reduction in water consumption through rainwater harvesting and closed-loop cooling systems \cite{Unilever2022}. Similarly, PT Pertamina's refineries have implemented advanced water treatment systems, achieving water recycling rates of 70-85\% across major facilities.

The pulp and paper sector has made significant progress in water efficiency. APRIL's Riau operations reduced water intensity by 25\% between 2015-2020 through process optimization and wastewater treatment improvements \cite{APRIL2021}. These efforts align with the Indonesia Water Mandate, a voluntary corporate water stewardship initiative launched in 2018.

\subsection{Regulatory Compliance and Standards}

Indonesian water management is governed by Law No. 17/2019 on Water Resources, which mandates permits for water withdrawal exceeding 5,000 cubic meters annually. The Ministry of Environment and Forestry (KLHK) enforces water quality standards through Government Regulation No. 22/2021, requiring industries to meet specific effluent parameters before discharge.

Companies must also comply with ISO 14001 environmental management standards, which include water usage monitoring and reduction targets. The Indonesian Sustainable Palm Oil (ISPO) certification requires oil palm plantations to implement water conservation measures and protect water catchment areas.

\subsection{Wastewater Treatment and Reuse}

Industrial wastewater treatment has improved significantly, with 85\% of large manufacturers operating on-site treatment facilities \cite{KLHK2022}. The Jakarta Industrial Estate Pulogadung reports 95\% wastewater treatment compliance among its tenants, with treated water either reused in operations or safely discharged to municipal systems.

Emerging technologies include membrane bioreactors and reverse osmosis systems, which enable water reuse rates exceeding 80\%. PT Freeport Indonesia's tailings management facility in Papua incorporates advanced water treatment, recovering 60\% of process water for reuse in mining operations.

\subsection{Climate Change Impacts and Adaptation}

Climate change poses significant risks to industrial water security in Indonesia. Changing rainfall patterns affect surface water availability, while sea-level rise threatens coastal industrial facilities through saltwater intrusion. The government's National Action Plan for Climate Change Adaptation (RAN-API) identifies water resource management as a priority sector.

Companies are increasingly incorporating climate resilience into water management strategies. PT PLN's power plants now include drought contingency plans, while beverage manufacturers like Coca-Cola Amatil Indonesia have diversified water sources and increased storage capacity to mitigate supply risks.

\subsection{Stakeholder Engagement and Transparency}

Corporate water disclosure has improved, with 68\% of Indonesia's largest companies reporting water metrics in sustainability reports \cite{ACCCSR2022}. The Indonesia Stock Exchange's ESG reporting guidelines recommend water usage disclosure, though adoption remains voluntary.

Community engagement around water resources has become increasingly important. PT Agincourt Resources' Martabe gold mine in North Sumatra established a community water committee to address local water access concerns, demonstrating how industrial water management intersects with social license to operate.

\subsection{Future Challenges and Opportunities}

Indonesia faces several water management challenges, including aging infrastructure, competing water demands, and limited enforcement capacity. The government's National Medium-Term Development Plan (RPJMN) 2020-2024 prioritizes water infrastructure investment, but implementation gaps remain.

Opportunities exist for circular water economy approaches, where industrial wastewater becomes a resource for other uses. The Surabaya Industrial Estate Rungkut successfully implemented industrial symbiosis, with treated wastewater from food processing plants used for cooling in nearby facilities.

Digital water management technologies offer additional efficiency gains. Smart metering and real-time monitoring systems can identify leaks and optimize usage, though adoption rates remain low due to cost barriers for smaller enterprises.

The integration of water stewardship into corporate strategy represents a maturing approach to water management. Companies increasingly recognize that sustainable water use supports both environmental objectives and operational resilience, particularly in water-stressed regions of Indonesia.